% 07_choosing_methods.tex
\section{Choosing and combining methods}

Run at least two of the three methods on every question, and treat close agreement as the main sign that a result is robust. Where they disagree, look first at which confounders each method adjusts for, and at overlap.

\begin{xltabular}{\textwidth}{YL{3.6cm}Y}
\toprule
\textbf{Question} & \textbf{Best suited} & \textbf{Why} \\
\midrule
Does the treatment work, explained to clinicians? & Target trial emulation & Transparent pairs and balance tables \\
How does benefit change with time to treatment or dose? & S-learner & Treatment enters as a continuous feature \\
Which patients benefit most? & Causal forest (with BLP test) & Designed for heterogeneity, with a check that it is real \\
Many confounders, some only weakly related to outcome & Causal forest or S-learner & Matching cannot include many variables exactly \\
Small treated group & Matching (with seed variation) or causal forest & Report seed variation; S-learner may shrink small-arm effects \\
Strong confounder recorded for only some patients & All three, as a sensitivity analysis & Restrict, or add with missing as its own category \\
\bottomrule
\end{xltabular}

\subsection{Recommended workflow}

\begin{enumerate}
\item Write the target trial protocol and draw a DAG; decide the role of each variable, especially timing variables.
\item Define the cohort so that untreated patients could genuinely have been treated (same eligibility window and time zero).
\item Run diagnostics before reading effects: matching rate and balance; model discrimination and calibration; overlap.
\item Estimate the ATT with at least two methods; report absolute and relative effects with honest CIs (paired, bootstrap or influence-function).
\item Run the sensitivity analyses listed for each method, changing one choice at a time, and the three shared confounding tools.
\item Compare methods in one table, and explain any disagreement.
\item Report using the TARGET statement \cite{cashin2025} for the emulation and TRIPOD \cite{collins2015tripod} for model performance.
\end{enumerate}

\subsection{Practical conventions}

\begin{itemize}
\item Keep every analysis choice (cut-offs, outcome definition, seeds, thresholds) in one settings block, so each sensitivity analysis changes one line.
\item Fix random seeds, and report variation across seeds rather than one seed.
\item Save every table to CSV as it is produced; long refits should save as they go.
\item Describe methods before each code block and results after, so notebooks double as the methods record.
\end{itemize}
